% =========================================================
% doc1-intro.tex — Section 1: Introduction
% Source: d0 §Sec1; d1 §1 (motivation, contributions, related work)
% =========================================================

\section{Introduction}\label{sec:intro}

Order-reversing involutions arise in two largely separate traditions. In the
theory of residuated structures, of which bounded Girard algebras are the
representative example, a structure carries a negation $\dstar{(\cdot)}$
obtained from the residual by fixing a distinguished element $\ep$ and setting
$\dstar{a} := a \res \ep$ \citep{Rosenthal1990,Agliano2025}. This negation has
two properties: it is an involution, $\dstar{(\dstar{a})} = a$, and it is
antitone with respect to the order, $a \leq b$ implying
$\dstar{b} \leq \dstar{a}$. In this tradition the residual is the source of
both properties: the involution and its order-reversal are consequences of the
residuation law, not primitive assumptions.

In algebra, by contrast, a semiring with involution is equipped with a negation
$\dstar{(\cdot)}$ as a primitive operation satisfying only
$\dstar{(\dstar{a})} = a$, that is, an anti-automorphism of the semiring
\citep{Golan1999,dolinka2000minimal}. Here the negation is primitive rather than
manufactured, but it is not antitone; indeed the structure carries no order at
all with respect to which antitonicity could be stated.

These two traditions leave a natural question open: is there a structure in
which the negation is a primitive operation possessing both properties at
once---involution and antitonicity---without being generated from a residual?

This paper answers the question affirmatively. We retain the involution axiom
of the algebraic tradition and replace the residual by a single new axiom
relating the negation to the join. On this basis the order is not primitive but
derived: it is recovered from the join, and antitonicity is not assumed but
follows. Concretely, we introduce \emph{Polar Semirings}: idempotent
commutative semirings equipped with a primitive negation subject to two axioms,
the involution axiom and a new polarity--join compatibility axiom. On this
structure (i)~the order and the antitonicity of the negation are generated
rather than assumed; (ii)~no residual is used; and (iii)~in contrast to a
semiring with involution, the negation is not an anti-automorphism of a single
structure but an isomorphism onto a new dual idempotent commutative semiring.

This is not the only work to obtain a primitive, order-compatible polarity.
Dual pairs of semirings, developed to study fuzzy sets, provide such a polarity
as well \citep{movckovr2025unification}; that system, however, imposes
additional distributivity axioms tailored to its intended application. Our
concern is not a particular application but the axiomatic question of how little
is required: we add only the two axioms governing the negation on top of a
standard semiring, isolating the polarity and discarding the distributivity
constraints entirely. This isolation makes it possible (i)~to separate the
contribution of the polarity itself from any distributivity assumption; (ii)~to
identify the minimal axiomatic cost of a primitive, order-compatible polarity;
and (iii)~to obtain a structure that does not presuppose completeness or the
extra distributivity laws of dual pairs of semirings.

\paragraph{Contributions.}
The paper makes three contributions.
\begin{enumerate}
  \item \emph{Definition and minimality} (\Cref{sec:ps}). We axiomatize Polar
  Semirings by eleven axioms, exhibit canonical models, and prove that the
  axiom system is independent: no axiom is a consequence of the remaining
  ten.
  \item \emph{Architecture} (\Cref{sec:arch}). We show that the axioms
  decompose into three layers: an algebraic layer (the polarity-free
  semiring reduct), a reflection layer, in which the involution axiom alone
  generates the dual operations and all De Morgan laws at no axiomatic cost,
  and an integration layer, in which the single polarity axiom fuses the
  reflected order with the original one. The decomposition culminates in the
  Fundamental Structure Theorem: every Polar Semiring canonically induces a
  dual Polar Semiring, the polarity is an isomorphism onto the dual, and the
  two induced orders coincide.
  \item \emph{Positioning} (\Cref{sec:positioning}). We locate Polar
  Semirings relative to bounded Girard algebras. Every bounded Girard
  algebra, upon forgetting its residual, satisfies all eleven axioms; and
  there exists a Polar Semiring on the unit interval that cannot be extended
  to a bounded Girard algebra. Hence bounded Girard algebras form a proper
  subclass, and Polar Semirings are a strict, residuation-free
  generalization.
\end{enumerate}

\paragraph{Related structures.}
The order-theoretic reduct of a Polar Semiring is a bounded lattice with an
antitone involution satisfying the De Morgan laws, that is, a bounded
i-lattice in the sense of \citet{Kalman1958}; distributivity is not assumed,
so De Morgan algebras \citep{Rasiowa1974} and Boolean algebras arise from
Polar Semirings only after adding further axioms. On the multiplicative side,
semirings with involution \citep{Golan1999} retain the involution but drop
the order-compatibility axiom entirely. A brief orientation table comparing
these classes precedes the detailed comparison with bounded Girard algebras
in \Cref{sec:positioning}. Throughout the paper we maintain a strict
terminological distinction, elaborated in \Cref{sec:arch}: a structure is
called \emph{involutive} when it satisfies only the equation
$\dstar{(\dstar{a})}=a$, and \emph{polar} when the involution is additionally
compatible with the induced order, in Birkhoff's classical sense
\citep{Birkhoff1940}.

\paragraph{Outline.}
\Cref{sec:ps} defines Polar Semirings, presents basic models, and proves
independence of the axioms. \Cref{sec:arch} develops the three-layer
architectural analysis and proves the Fundamental Structure Theorem.
\Cref{sec:positioning} carries out the comparison with bounded Girard
algebras, proving both the inclusion and its strictness.
\Cref{sec:conclusion} summarizes and discusses future directions.
